%! AP Statistics — Unit 9, Lesson 9.6 — Group Activity (Answer Key)
\documentclass[10pt]{article}
\usepackage{apstats-article}
\usepackage{apstats-key}

\begin{document}

\pageheader{Unit 9, Lesson 9.6}{Group Activity --- Answer Key}
\namepartnerperiod

\vspace{0.1in}
\textbf{Inference Procedure Identification --- 8 Scenarios}

\textbf{Directions:} Read each scenario. Classify the data and identify the correct inference
procedure at your tier. Work with your group and use the decision framework from notes.

\vspace{0.1in}

\textbf{Scenarios:}
\begin{enumerate}[itemsep=4pt]
  \item A researcher records whether each of 80 randomly selected college students owns a car
        (yes/no). She wants to \textit{estimate} the proportion who own cars.
  \item A school surveys 120 students and records their preferred lunch option (pizza, salad,
        pasta, other). She wants to test whether the four options are equally popular.
  \item A researcher randomly assigns 40 patients to medication A or medication B and records
        recovery time (days) for each. She tests whether the two medications differ in mean
        recovery time.
  \item 50 pairs of identical twins are studied; one twin from each pair exercises 5+ hours per
        week and the other exercises less. The researcher records reaction time for each twin
        and tests whether exercise is associated with faster reaction time.
  \item A political scientist randomly surveys 300 voters in two cities (150 from each) and
        records party affiliation (Democrat, Republican, Independent). She tests whether the
        party affiliation distribution is the same in both cities.
  \item A nutritionist records daily sugar intake (g, $x$) and body fat percentage (\%, $y$)
        for a random sample of 35 adults. She wants to \textit{estimate} the true slope $\beta$.
  \item A researcher surveys 200 people and records education level (HS, Some College,
        Bachelor's, Graduate) and income category (Low, Middle, High) to test whether
        education and income are associated.
  \item A sports scientist randomly samples 60 high school athletes and records their resting
        heart rate (bpm). She wants to test whether the mean resting heart rate differs from
        the national average of 72 bpm.
\end{enumerate}

\vspace{0.1in}

\begin{tcolorbox}[colback=white, colframe=black!40,
    title=\textbf{Tier R --- Remediate},
    fonttitle=\bfseries, arc=2mm, left=3mm, right=3mm, top=2mm, bottom=2mm]

\renewcommand{\arraystretch}{2.0}
\begin{tabularx}{\linewidth}{@{} c X X X c @{}}
  \rowcolor{sky}
  \textbf{\#} & \textbf{Categorical or Quantitative?} & \textbf{Groups / Structure} & \textbf{Procedure Name} & \textbf{CI or Test?} \\
  \hline
  1 & \ans{Categorical} & \ans{1 proportion} & \ans{1-prop $z$-interval} & \ans{CI} \\
  2 & \ans{Categorical} & \ans{1 var, 4 categories} & \ans{Chi-sq GOF test} & \ans{Test} \\
  3 & \ans{Quantitative} & \ans{2 independent groups} & \ans{2-sample $t$-test} & \ans{Test} \\
  4 & \ans{Quantitative} & \ans{2 paired (matched)} & \ans{Paired $t$-test} & \ans{Test} \\
\end{tabularx}
\end{tcolorbox}

\begin{tcolorbox}[colback=white, colframe=black!40,
    title=\textbf{Tier A --- Approaching Proficiency},
    fonttitle=\bfseries, arc=2mm, left=3mm, right=3mm, top=2mm, bottom=2mm]

\renewcommand{\arraystretch}{2.0}
\begin{tabularx}{\linewidth}{@{} c X X X c @{}}
  \rowcolor{sky}
  \textbf{\#} & \textbf{Categorical or Quantitative?} & \textbf{Groups / Structure} & \textbf{Procedure Name} & \textbf{CI or Test?} \\
  \hline
  1 & \ans{Categorical} & \ans{1 proportion} & \ans{1-prop $z$-interval} & \ans{CI} \\
  2 & \ans{Categorical} & \ans{1 var, 4 categories} & \ans{Chi-sq GOF test} & \ans{Test} \\
  3 & \ans{Quantitative} & \ans{2 independent groups} & \ans{2-sample $t$-test} & \ans{Test} \\
  4 & \ans{Quantitative} & \ans{2 paired (matched)} & \ans{Paired $t$-test} & \ans{Test} \\
  5 & \ans{Categorical} & \ans{1 var, 2 populations} & \ans{Chi-sq homogeneity} & \ans{Test} \\
  6 & \ans{Quantitative} & \ans{Bivariate / slope} & \ans{$t$-interval for slope} & \ans{CI} \\
  7 & \ans{Categorical} & \ans{2 categorical vars, 1 sample} & \ans{Chi-sq independence} & \ans{Test} \\
  8 & \ans{Quantitative} & \ans{1 mean} & \ans{1-sample $t$-test} & \ans{Test} \\
\end{tabularx}

\vspace{0.1in}
\textbf{Scenario 1} --- \ans{Estimate $p$ = the true proportion of college students who own a car.}

\textbf{Scenario 3} --- \ans{$H_0$: $\mu_A = \mu_B$ (no difference in mean recovery time); $H_a$: $\mu_A \ne \mu_B$.}

\textbf{Scenario 6} --- \ans{Estimate $\beta$ = the true slope of the population regression line relating sugar intake to body fat percentage.}
\end{tcolorbox}

\begin{tcolorbox}[colback=white, colframe=black!40,
    title=\textbf{Tier E --- Extension},
    fonttitle=\bfseries, arc=2mm, left=3mm, right=3mm, top=2mm, bottom=2mm]

Complete all of Tier A above. Then answer the following.

\vspace{0.08in}
\textbf{Bonus 1.} \ans{The twins are naturally paired by genetics. Because each pair shares the same genetic background, the two observations in each pair are not independent. A paired $t$-test uses the \textit{differences} within each pair as the data, removing the variability due to genetics and increasing statistical power. A 2-sample $t$-test would ignore this pairing and be less powerful.}

\textbf{Bonus 2.} \ans{A chi-square test for \textit{homogeneity} is used when data come from two or more separate populations (or groups) and we want to compare the distribution of a single categorical variable across those groups (Scenario 5: two cities). A chi-square test for \textit{independence} is used when data come from one sample and we want to determine whether two categorical variables are associated (Scenario 7: one sample, two variables recorded).}
\end{tcolorbox}

\end{document}
